\begin{abstract}
\noindent Per- and polyfluoroalkyl substances (PFAS) in surface waters have entered a regime of enforceable numeric limits---a 2024 U.S.\ drinking-water Maximum Contaminant Level of 4~ng~L\textsuperscript{-1} for PFOA and PFOS, a CERCLA hazardous-substance designation, and state in-stream criteria---shifting the management question from detection toward reach-scale prediction and source attribution, tasks that monitoring and statistical occurrence models structurally cannot supply. We present a process-based PFAS surface-water fate-and-transport engine implemented in modern (free-form) SWAT\textsuperscript{+}. The engine solves a soil three-phase equilibrium---aqueous, Freundlich solid-phase, and Langmuir air--water interfacial adsorption---for each hydrologic response unit and soil layer, mobilizes PFAS through surface runoff, lateral flow, leaching, and sediment-bound pathways, and routes it through the channel network with a linear organic-carbon partition and benthic settling, resuspension, diffusion, and burial; reservoirs are treated as mass-conserving mixed cells and point sources (e.g.\ contaminated-site leachate, effluent) are injected per reach. Two correctness checks anchor the implementation: the equilibrium solver agrees with a 40-digit reference to a worst-case relative error of $1.7\times10^{-7}$ across 29{,}160 cases, and the in-stream routine is numerically identical to the established SWAT\textsuperscript{+} pesticide channel solver---restricted to its linear-partition subset---to single-precision floating-point tolerance, so the new physics is the soil column rather than the channel mathematics. We demonstrate the engine on a high-resolution SWAT\textsuperscript{+} model of the Rogue River watershed (Michigan; USGS~04118500), generated automatically from NHDPlus High Resolution hydrography and national environmental layers and resolving the network into 590 channels and 17{,}771 hydrologic response units, which lets each of 31 ambient PFOS observations be assigned to its hydrologically correct reach. With a single global soil-loading scale factor the model reproduces the observed Rogue mainstem PFOS gradient---rising downstream toward the Wolverine Worldwide contaminated-site corridor---at a log-space root-mean-square error of $0.15$~dex along the seven source-bearing mainstem reaches ($0.74$~dex across all $20$ gauged reaches, $29$ stations), while the residual underprediction at the source corridor and the scatter among tributaries are physically interpretable as, respectively, the concentrated point source and per-land-use soil heterogeneity. A $120$-member Latin-hypercube uncertainty ensemble confirms this decomposition: the predictive envelope encloses five of the seven mainstem observations ($71\%$) but almost none on the tributaries, locating the unexplained variance in local soil heterogeneity rather than in the global parameters, of which soil loading dominates. The contribution is a transferable, mass-conserving causal layer for PFAS that complements monitoring and machine-learning triage; reported skill is a first watershed demonstration rather than a national validation, and expert review is appropriate before decision use.
\end{abstract}
